\setcounter{aufgabe}{13}
\einheitenkopf{Einheit 2 von 5 · Säulen-, Balken- und Liniendiagramme}

\begin{aufgabe}{Diagramm -- Was ist ein Kästchen wert? Gib nur an, wie viel ein Schritt von einer Linie zur nächsten wert ist. Lies keinen Wert ab.}
\begin{teile}
\teil Die Linien sind mit $0$, $10$, $20$, $30$ beschriftet, dazwischen liegt keine weitere Linie. \leerfeld
\teil Die Linien sind mit $0$, $100$, $200$ beschriftet, dazwischen liegt je eine Hilfslinie. \leerfeld
\teil Die Linien sind mit $0$, $20$, $40$, $60$ beschriftet, dazwischen liegen je drei Hilfslinien. \leerfeld
\teil Die Linien sind mit $0$; $0{,}05$; $0{,}10$ beschriftet, dazwischen liegen je vier Hilfslinien. \leerfeld
\teil Die Linien sind mit $0$, $200$, $400$ (in Mio.) beschriftet, dazwischen liegen je drei Hilfslinien. \leerfeld
\end{teile}
\end{aufgabe}

\begin{aufgabe}{Diagramm -- Wo fängt die Achse an? Kreuze an, ob die senkrechte Achse bei null beginnt. Wenn nicht, schreibe den Startwert auf.}
\begin{teile}
\teil Beschriftung $0$, $5$, $10$, $15$ \hfill \janein \quad Start: \leerfeld
\teil Beschriftung $40$, $45$, $50$, $55$ \hfill \janein \quad Start: \leerfeld
\teil Beschriftung $0$, $250$, $500$ \hfill \janein \quad Start: \leerfeld
\teil Beschriftung $3{,}10$; $3{,}15$; $3{,}20$ \hfill \janein \quad Start: \leerfeld
\teil Beschriftung $820$, $830$, $840$ \hfill \janein \quad Start: \leerfeld
\end{teile}
\end{aufgabe}

\begin{aufgabe}{Säulendiagramm -- ablesen. Das Diagramm zeigt, wie viele Gäste eine Eisdiele an den Tagen einer Woche hatte.}
\noindent\saeulenab[ymax=100,ystep=20,yfein=5,ylabel=Gäste]{Mo/20,Di/60,Mi/40,Do/80,Fr/55,Sa/75,So/65}

\begin{teilezwei}
\tz Montag \leerfeld & \tz Dienstag \leerfeld \\
\tz Mittwoch \leerfeld & \tz Donnerstag \leerfeld \\
\tz Freitag \leerfeld & \tz Sonntag \leerfeld \\
\end{teilezwei}
\begin{teile}
\teil An wie vielen Tagen kamen \emph{mehr als} 60 Gäste? \leerfeld
\teil An welchem Tag kamen die meisten, an welchem die wenigsten Gäste?
\par meiste: \leerfeld \quad wenigste: \leerfeld
\teil Wie viele Gäste kamen am Samstag mehr als am Mittwoch? \leerfeld
\end{teile}

Das zweite Diagramm zeigt die Zuschauerzahlen eines Stadions. Achte auf die Beschriftung der Achse.

\noindent\saeulenab[ymax=50,ystep=10,yfein=2,ylabel=Zuschauer in Tausend,breite=1.6]{2021/30,2022/24,2023/38,2024/46}

\begin{teile}
\teil Wie viele Zuschauer kamen 2022? \leerfeld
\teil Wie viele Zuschauer kamen 2024? \leerfeld
\teil Wie viele Zuschauer kamen 2024 mehr als 2021? \leerfeld
\end{teile}
\end{aufgabe}

\begin{aufgabe}{Balken- und Liniendiagramm -- ablesen. Das Balkendiagramm zeigt, wie oft vier Farben als Lieblingsfarbe genannt wurden.}
\noindent\balkenab[xmax=40,xstep=10,xfein=2,xlabel=Nennungen]{Rot/24,Blau/32,Grün/14,Gelb/8}

\begin{teilezwei}
\tz Rot \leerfeld & \tz Grün \leerfeld \\
\end{teilezwei}
\begin{teile}
\teil Welche Farbe wurde am häufigsten genannt? \leerfeld
\end{teile}

Das Liniendiagramm zeigt die mittlere Temperatur der ersten sechs Monate eines Jahres.

\noindent\liniendia[ymax=20,ystep=5,yfein=1,ylabel=Temperatur in $^\circ$C,breite=1.4]{Jan/2,Feb/5,Mär/4,Apr/9,Mai/14,Jun/17}

\begin{teile}
\teil Wie warm war es im März? \leerfeld
\teil In welchem Monat war es am kältesten? \leerfeld
\teil Zwischen welchen beiden Monaten ging die Temperatur zurück? \leerfeld
\teil Beschreibe in einem Satz, wie sich die Temperatur von März bis Juni verändert. \\[10pt]
\end{teile}
\end{aufgabe}

\begin{aufgabe}{Diagramm -- zeichnen. Ein Sportverein hat gezählt, wie viele Kinder die einzelnen Sportarten wählen.}
\sachtabelle{lccccc}{Sportart & Handball & Tennis & Judo & Reiten & Schwimmen}{%
Anzahl & 15 & 25 & 10 & 30 & 20}

Zeichne die Säulen in das vorbereitete Diagramm.

\noindent\saeulenab[ymax=40,ystep=10,yfein=5,ylabel=Anzahl,breite=1.7]{Handball/,Tennis/,Judo/,Reiten/,Schwimmen/}

\begin{teilezwei}
\tz Handball & \tz Tennis \\
\tz Judo & \tz Reiten \\
\end{teilezwei}
\begin{teile}
\teil Schwimmen
\teil Eine sechste Sportart wird von 75 Kindern gewählt. Was müsstest du an der senkrechten Achse ändern, damit auch diese Säule ins Diagramm passt? \\[10pt]
\end{teile}
\end{aufgabe}

\begin{aufgabe}{Diagramm -- Achseneinteilung. An der senkrechten Achse fehlen die Zahlen. Bekannt ist: Am Montag kamen 120 Besucher.}
\noindent\saeulenab[ymax=240,ystep=30,yfein=30,ylabel=Besucher,ohnezahlen,breite=1.6]{Mo/120,Di/180,Mi/90,Do/}

\begin{teile}
\teil Wie viele Besucher zeigt ein Kästchen? \leerfeld
\teil Wie viele Besucher kamen am Dienstag? \leerfeld
\teil Wie viele Besucher kamen am Mittwoch? \leerfeld
\teil Am Donnerstag kamen 210 Besucher. Zeichne die Säule ein.
\teil Das Diagramm unten zeigt für vier Sprachen, wie viele Menschen sie als Muttersprache sprechen. Beschriftet ist nur die Säule für Spanisch: 480 Mio. Außerdem ist bekannt: Hindi sprechen 320 Mio., Portugiesisch 240 Mio. und Arabisch 400 Mio. Menschen. Bestimme, wie viele Mio. ein Kästchen zeigt, schreibe „Hindi“ und „Portugiesisch“ anstelle der beiden Fragezeichen und zeichne die Säule für Arabisch ein. (P10 2023)
\end{teile}

\noindent\saeulenab[ymax=560,ystep=80,yfein=80,ylabel=Mio. Menschen,ohnezahlen,breite=1.8]{Spanisch/480,?/320,?/240,Arabisch/}
\end{aufgabe}

\begin{aufgabe}{Diagramm -- Fehler finden.}
An der senkrechten Achse eines Diagramms steht „Zuschauer in Tausend“. Die Säule für Samstag reicht genau bis zur Linie 18. Jonas notiert:
\rechnung{\text{Am Samstag kamen 18 Zuschauer.}}
\begin{teile}
\teil Finde den Fehler, benenne ihn und gib die richtige Zahl an. \\[10pt]
\teil Die Säule für Sonntag reicht auf derselben Achse bis 25. Wie viele Zuschauer kamen am Sonntag? \leerfeld
\end{teile}
\end{aufgabe}

\begin{aufgabe}{Diagramm -- Begründen.}
\begin{teile}
\teil Die Zahl der Geburten in einer Stadt soll für die Jahre 2015 bis 2024 dargestellt werden. Begründe, warum sich dafür ein Säulen- oder Liniendiagramm besser eignet als ein Kreisdiagramm. \\[12pt]
\teil Zwischen den beschrifteten Linien $0$ und $50$ liegen vier Hilfslinien. Jemand sagt: „Dann ist eine Hilfslinie 5 wert.“ Entscheide ohne Rechnung, ob das stimmen kann, und begründe. \\[12pt]
\end{teile}
\end{aufgabe}
